% Chapter 10: Towards deployment: severity and inference cost.
% Sources: results/severity/SEVERITY_SUMMARY.md (57e2bb4c), thesis/INFERENCE_COST.md (1e5a9002), tables T18, T22, T9,
% benchmark/release_v1/LICENCE_AUDIT.md (f267ea23), PROTOCOL_v1.0.md (eeb3379a) s2.4. Every number: ch10_deployment_NOTES.md.
% Bibliography: ch10_deployment.bib, plus deitsch2019elpv and xie2021segformer from ch06_protocol.bib and wang2021hrnet from refs_extra.bib.

\chapter{Towards deployment: severity and inference cost}
\label{ch:deployment}

A model that the industry partner could use for inspection needs more than a crack score: an output a person can act on, a known cost per cell,
and terms that allow the use. This chapter reports what we measured on each. The severity work is descriptive (track T5 of the
protocol): it used the PF baseline out of fold on pool cells and three external data sets for inference and analysis only, and
trained no model. The cost numbers are those of the E1 candidates of Chapter~8. Tables~\ref{tab:T18} and~\ref{tab:T22} hold the
values.

\section{The severity output}
\label{sec:deploy-output}

The segmentation gives, per cell, an inactive area fraction $a$: the predicted inactive area, plus the crack pixels it encloses, as a
percentage of the cell, with a band (0, under 1\%, 1\% or more) and an error bar from our out-of-fold record. Under the protocol's
scoring unit, ``cell'' here means the image window (Chapter~6). Per module, the output is the worst cell (largest $a$) with its band and the module sum of $a$ in cell equivalents. An isolation-risk
score (cells cut off from the ribbons by cracks) was dropped: it works on clean labels but not on our predictions. No power band or
watt figure is given, because no labelled cell of ours has an electrical measurement.

On the 532 pool cells, with the PF baseline out of fold (Figure~\ref{fig:F18}), the mean absolute error of $a$ is 0.12 percentage points, with a bias of
$-$0.08 points (the model understates $a$). Most cells have no inactive area; on the 48 cells that do, the error is about 1.25 points
(a post-hoc split). Band agreement is kappa 0.78. At a detection threshold of 0.5\%, sensitivity is 0.73 at a false-alarm rate of
0.25\%, but on CSIR cells sensitivity is only 0.20. The ``under 1\%'' band is unreliable; cells predicted at 1\% or more are truly at
1\% or more in 93\% of cases. These are pool cells as built: the output is not validated on unseen modules, on half-cut cells or on
the final model C2, for which no severity number exists. We note, as judgement, that E1 is 101 of 106 CSIR cells, so the source on
which the output is weakest is the main source of E1.

\FIG{F18}{Severity output of the PF baseline. (a) Predicted against labelled inactive area fraction $a$ per window, 532 pool cells; PF
out of fold on cv5v2 (not cv5v3), mean over seeds 5 to 9; symlog axes, dashed line $y=x$; 467 cells sit at (0, 0). (b, c) ELPVPower
modules (129 rows used): relative power loss against predicted inactive area and against predicted crack length (PF seed 5, 5-fold
ensemble); pre-registered outcome NULL. Not validated on unseen modules, half-cut cells or the final model.}

\section{Does the output say anything about power?}
\label{sec:deploy-power}

The pre-set question was whether predicted inactive area predicts power loss better than predicted crack length. Three public sets
with electrical measurements were used, for inference and analysis only.

\paragraph{Polenghi.} Two small polycrystalline panels pass through 84 cumulative damage stages, with power at one operating point
under a lamp of about 50~W/m$^2$ \citep{polenghi2026data}. Predicted inactive area tracks the power loss at Spearman $\rho$ 0.56,
predicted crack length at 0.65, and the stage number alone at 0.91. With $n=2$ devices and damage that only accumulates, this record
cannot tell damage measures apart: NULL. The model also under-calls large dark strips here, so the null is not evidence against
inactive area.

\paragraph{ELPVPower.} This set holds 137 module instances with measured maximum power \citep{hoffmann2020elpvpower}; it was used for
analysis only, and its CC~BY licence version is not stated. Predicted inactive area ranks modules by power loss at $\rho$ 0.87
[0.80, 0.91] (129 modules); predicted crack length does as well (0.89), and the order of the three module types alone gives 0.83.
Within the near-clean type, nothing the model sees tracks the loss ($\rho$ 0.02). A look within the heavily cracked type (59
modules) favours inactive area (0.93 against 0.83), but it was taken after the fact and is not a result. Under the pre-set rule the
comparison is NULL. A duplicate check found no ELPVPower cell among our training cells.

\paragraph{DTU A3.} This set holds lab-damaged modules with STC I-V curves \citep{delpradosantamaria2025drone}; the label is the
fill-factor (FF) loss, which misses bypass-diode loss. A simple count of EL darkness tracks FF loss on the 24 non-diode modules at
$\rho$ 0.64 [0.24, 0.87]. On the 16 modules our cell cutting could handle, predicted inactive area tracks FF loss at $\rho$ 0.90,
against 0.44 for EL darkness, which meets the pre-set bar by a hair. But the interval's lower end is about 0 and depends on the
bootstrap seed, the four highest-loss interconnect modules (where EL darkness works best) failed cutting, and 9 of the 16 modules have
FF losses within about 0.8\% of zero. We read it as borderline, on a selected subset.

Taken together, ``predicted inactive area predicts power better than crack length'' is not supported: null on Polenghi, null on
ELPVPower, borderline on a subset of A3 (Table~\ref{tab:T20}, claim 21). Turning $a$ into watts or a power band, rating a module from
our labelled crops, and anything on half-cut modules (cell cutting failed on them in A3) are not supported either.

\section{What thermal imaging adds}
\label{sec:deploy-thermal}

The industry partner owns a drone thermal camera, so we asked what drone thermal adds, at module level, on A3. The hot-spot spread (99th
percentile minus median temperature) does not track FF loss for cracks, interconnect faults or potential-induced degradation ($\rho$
0.19), and added to EL it shows no clear gain ($\Delta R^2$ +0.10, interval $-$0.18 to +0.40). It does flag bypass-diode faults and
soiling strongly: 16 to 38\,$^\circ$C, against about 1\,$^\circ$C for the other defects. Thermal is therefore a screen for diode and
soiling faults, which EL-based severity does not cover; it is not a crack-severity measure. We make no cell-level thermal claim.

\section{Inference cost}
\label{sec:deploy-cost}

We timed the network part of the E1 candidates: the whole prediction for one batch (all members and views, the D4 transforms and
their inverses, softmax and averaging) divided by the batch size, without image reading, preprocessing, argmax or writing. Inputs were
random tensors of the real shape ($1\times512\times512$), so no image or label was read. The CPU runs used one AMD Threadripper 7960X
(24 cores) at 16 threads and batch 1; the GPU runs used one RTX A6000 in fp32 with TF32 off, as for the sealed predictions, at batch 1
and 8, with 200 timed cells per configuration. Both hosts were shared; host load can only slow a run, so the times are upper bounds
for idle hardware at these settings (Table~\ref{tab:T22}).

PF has 26.1~M parameters and costs 148.0 GFLOPs per view at 512~px; HRNet-W18 U-Net++ has 20.8~M and costs 218.1. Per cell, C1 costs
148 GFLOPs, the PF ensemble E-PF 740, C1-D4 1,184, C2 1,744 and the C2 ensemble C3 8,722. C2 is $8\times$ C1 in forward passes and
$11.8\times$ in FLOPs. On the CPU the median network time per cell is 0.216~s for C1, 2.23~s for C1-D4 and 3.88~s for C2 ($18.0\times$
C1). On the GPU at batch 1 it is 0.0112~s for C1, 0.0911~s for C1-D4 and 0.174~s for C2 ($15.5\times$ C1); at batch 8, C2 takes 0.133~s per cell, and batching its eight views into one forward pass gives 0.134~s. One HRNet view costs
$2.3\times$ a PF view on the GPU and $1.9\times$ on the CPU, against $1.47\times$ in FLOPs, so FLOPs understate HRNet's time on both
devices. C2's peak GPU allocation is 445~MiB at batch 1 and 2,993~MiB at batch 8; its peak resident memory on the CPU at batch 1 is
1,376~MiB. At batch 1, one GPU runs C2 on about 20,700 cells per hour and one CPU on about 930.

Network time is not the whole job. The sealed E1 jobs of one C2 seed took about 100~s for 106 cells, including start-up, checkpoint
and data loading and the writing of log-probabilities; network time ($106\times0.174=18$~s) is under a fifth of that.

\FIG{F19}{Cost against quality for the E1 candidates: (a) GFLOPs and (b) GPU time per cell (one RTX A6000, fp32, batch 1, median),
both on log axes, against E1 crack relF1@2. C1, C1-D4 and C2 are means over seeds 5 to 9, with the range over seeds as a vertical line
(not an interval); E-PF and C3 are one ensemble each. Levels of different models on one set; no interval for the differences beyond
those of Table~\ref{tab:T9}. C2 recipe averaged over seeds 5 to 9; OP-A; 101 of 106 cells CSIR.}

What the cost buys (Figure~\ref{fig:F19}) is limited by what E1 can resolve. On E1 (unseen modules, supported, not proven; recipe averaged over seeds 5 to 9;
OP-A only; 101 of 106 cells CSIR; carried by the full cells) C2 beats C1 by +0.0430 [+0.0013, +0.1043] crack relF1@2, CONFIRMED
(Chapter~8). Most of the gain is the model: C2 $-$ C1-D4, both with eight views, is +0.0451 [+0.0173, +0.0699]. The D4 part is not
resolved on E1: C1-D4 $-$ C1 is $-$0.0021 [$-$0.0423, +0.0575]. The ensemble C3 (40 forward passes and 8,722 GFLOPs per cell, against
C2's 8 and 1,744) adds nothing measurable over its members ($-$0.0025 [$-$0.0254, +0.0157]). One checkpoint, C2 seed 5, is not
separated from C1 (+0.0189 [$-$0.0403, +0.1068]): the confirmed claim is about the recipe, not one checkpoint. Our judgement: since the
ensemble adds nothing, the deployable unit is one C2 checkpoint (seed 5, designated before any number), at about 0.17~s of network
time per cell on one GPU, and its expected gain rests on the recipe result. Whether the eight D4 views, which make up the whole
$8\times$ forward count, earn their cost on new modules is not resolved by E1; on CV5 they add +0.0055 [$-$0.0002, +0.0115] to the
C2 recipe's fold models (NOT CONFIRMED, NO LOSS; selection-biased).

\section{Licence: what the industry partner may deploy}
\label{sec:deploy-licence}

As trained, C2 is a research-only model. Its training pool contains the 50 BMRK cells, which are ELPV images \citep{deitsch2019elpv}
licensed CC BY-NC-SA 4.0; commercial use needs the holder's permission. The licence audit judges that this likely extends to models
trained on those images, a point untested in law. BenchmarkEL, which holds all pool cells, has no licence file and asks only for a
citation; the E-SCDD corrections are MIT on their dataset card, but that claim cannot clear the BenchmarkEL images. C2's pretrained
HRNet-W18 weights \citep{wang2021hrnet} are MIT. The SegFormer MiT-B2 weights \citep{xie2021segformer}, used only in one benchmark row
and not in C2, are licensed for non-commercial research and evaluation only.

The decision of 22 November is that no commercial model is trained on the BMRK cells or uses MiT-B2 weights. A model the industry partner could
sell therefore needs the C2 recipe retrained without the 50 BMRK cells, or the holder's permission; BenchmarkEL's own terms for
commercial use stay unclear until its authors answer. The effect of dropping the BMRK cells was not measured. A retrained model is
a new model: CV5 can describe it, but E1 is no longer a zero-shot set for us after opening \#9 (Chapter~6), so its unseen-module quality would need new labelled
data.

\section{What a pilot with the industry partner must measure}
\label{sec:deploy-pilot}

The external sets show what a pilot on the industry partner's own modules would need (judgement from the numbers above). Per module: STC I-V
(maximum power, short-circuit current, open-circuit voltage, fill factor) with module type and nameplate power; EL at a known
current, at two currents if possible; drone thermal at 600~W/m$^2$ or more; and a repeat I-V on a few modules. The modules should vary in cracks and inactive area independently within one
module type, since pooling types let type drive the ELPVPower result, and should include interconnect faults and half-cut modules,
with a cell-cutting tool that handles them. About 31 modules are needed for a true $\rho$ of 0.5 to give an interval above 0, and
about 90 for 0.3. The rule should be written before measuring, with one primary: the rank correlation of the module sum of predicted
inactive area with relative maximum-power loss, within type, against crack length. As of 4 November the industry partner had drone thermal only,
with an EL rig 2 to 3 months away; until it arrives a pilot can test the thermal screen, not EL severity. The segmentation side has a
matching gap: on the 62 half-cut cells of E1 the gain of C2 over C1 is not detected (+0.0110 [$-$0.0289, +0.0722]), so a pilot on
half-cut modules would also test the model.
